% =====================================================================
% Discussion & limitations (Sec 9) - material sources:
%   * D1 review sec 1.4 item 9 (four objection classes), item 6
%     (metatheorem gap), item 8 (reproducibility package)
%   * p0_reconciliation_report.md sec 4 (five-class decomposition,
%     71.4% no-content headline; review-risk note)
%   * sorry_contamination_audit.md (30.2% reversal, full-population)
%   * proof_status.md sec 5 (proof-rate caliber history)
%   * framework/journal_matching_matrix.md (D2: CPP 2028 main,
%     ITP 2027 mid-way; JAR main target, LMCS backup)
%   * new_directions_roadmap.md sec D (W-line future program)
% Number policy: all decompositions canonical per P0; single-case
% audit-population figures quoted with their sources.
% =====================================================================

\subsection{Caliber discipline as the central thesis}
\label{sec:thesis}
The paper's transferable claim is not the specific numbers of one
library but a methodological proposition: in a formal library
that mixes human and machine authorship, a proof-progress metric
is meaningful only as a (number, caliber) pair, with the caliber
declared, machine-checkable, and enforced by the same CI that
enforces the rest of the honesty layer. The program's two
self-caught incidents are the evidence.

The 30.2\% ``sorry contamination'' alarm~\cite{sylvasorryaudit2026}
was produced by a
string-level scan over the batch corpus: 36{,}199 of 119{,}859
files ``contained sorry''. A full-population re-audit under the
code-level caliber found the true count to be zero: 36{,}198 of
the 36{,}199 hits were the fixed header comment \emph{asserting
the absence} of placeholders, and the one remaining hit was a
theorem whose \emph{name} contains the word. The alarm had been
measuring comments, not proofs --- and the operational consequence
cuts both ways: a filter built on the string caliber would have
quarantined roughly a third of a clean corpus. The dashboard
history tells the same story in percentages: the headline rate
climbed 88.08\% $\to$ 99.5\% $\to$ 100\% across successive
documents before the honest re-definition dropped it to 3.91\%.
Every intermediate value was defensible under its own (undeclared)
caliber; the sequence is absurd only against a fixed one.

The redemption ledger closes the argument. Decomposed into the
five-class taxonomy~\cite{sylvap0recon2026}, the 196 closures are: 50 vacuous placeholder
closures (25.5\%), 90 $P\!\to\!P$ conditionalizations (45.9\%),
17 data-to-definition conversions (8.7\%), 24 definitional
reconstructions (12.2\%), and 13 substantive proofs (6.6\%), plus
2 bundling restructurings excluded from the content count. Each of
the two large classes is individually defensible --- a vacuous
placeholder was never mathematics, and a $P\!\to\!P$
conditionalization is a conservative extension that at least
makes the dependence explicit in the statement's type --- yet
together they are proxy-clean, and 140/196 (71.4\%) of closures
carry no proof content at all. A redemption program optimized on
any aggregate metric would drift exactly here: conditionalization
removes an \texttt{axiom} declaration, improves every census, and
transfers nothing. A skeptic who suspects a proof-rate metric of
being gamed now has a concrete audit program: decompose the
numerator into semantic classes, declare the caliber of every
count, and reconcile ledger, process, and live-code totals. That
is the paper's export.

\subsection{Limitations}
\label{sec:limitations}
This section states, rather than defends, what the program has
not achieved.

\begin{enumerate}
\item \textbf{Most closures transfer no mathematics.} The
canonical decomposition is unflattering by design: 140 of 196
redemptions (71.4\% --- the 90 $P\!\to\!P$ conditionalizations and
50 vacuous closures) add no proof content, and the 13 substantive
proofs (6.6\%) are themselves shallow --- arithmetic rewrites,
existence witnesses, and single-lemma mathlib discharges. The
declaration census improving 445$\to$253 is a record of
bookkeeping progress, not of mathematical progress, and this
paper deliberately refuses the aggregate ``196 redeemed''
reading.
\item \textbf{$P\!\to\!P$ conditionalization is a loophole the
taxonomy only exposes.} The conversion \texttt{axiom X : P} $\to$
\texttt{theorem X (h:P) : P := h} is logically conservative and
audit-proof: it removes an axiom, changes no theorem, and cannot
be faulted under any declaration-level metric. A library could
convert its entire axiom debt this way and truthfully report
``zero axioms'' while having verified nothing new. The five-class
taxonomy makes such a program \emph{visible}; it does not yet
\emph{prevent} it --- the CI gate audits the honesty of reporting,
not the depth of content.
\item \textbf{The conservativity claims are not yet metatheorems.}
That the three conversion families (conditionalization, vacuous
closure, definitional reconstruction) are conservative extensions
is asserted from inspection of the sweep records, not proved.
Until the small metatheory is formalized (the registered P1-6
program), the taxonomy's semantics rest on exactly the kind of
undischarged trust this paper is about retiring.
\item \textbf{The residual debt is large and partly permanent.}
253 axioms remain (239 primitive, 14 schema), including the
spectral-bridge axiom of Case Study I, which is physics rather
than mathematics and may never be ``redeemed'' in the sweep sense.
The Agda track's 131 postulates have never compiled in CI (the
toolchain OOMs on the audit machines). The honest proof rate of
3.91\% is itself the statement that the corpus is at an early
stage of genuine coverage.
\item \textbf{Verification boundaries.} The honesty CI checks the
honesty layer (markers, ledgers, manifests, caliber artifacts);
it does not run \texttt{lake build}. Sweep6 executed without a
Lean toolchain (structural verification only); the T3$'$/T6/T7
strengthening of Case Study I awaits its independent compile.
Reproduction is command-level and version-pinned, not yet
containerized (P1-8).
\item \textbf{Single corpus, no baseline benchmark.} The
methodology is evaluated on one library --- this one. The
comparison against naive \texttt{grep} and mathlib's linters is
qualitative (Sec.~\ref{sec:related}); a measured
detection/false-positive baseline comparison is future work. All
figures are anchored to the v7.99 working tree and will drift;
the caliber discipline, not the numbers, is what transfers.
\end{enumerate}

\subsection{Anticipated objections}
\label{sec:objections}
Four objections the internal review pre-registered~\cite{sylvad1review2026},
with the program's standing answers:

\begin{enumerate}
\item \emph{``The corpus is AI-generated; why trust any of it?''}
That is the premise, not the objection. The governance layer
exists because AI-assisted generation at this scale cannot be
trusted through review capacity. Trust is replaced by
verifiability: every headline number is a command away
(Appendix~\ref{sec:repro}), every claim row links its evidence,
and every artifact carries its provenance markers.
\item \emph{``The remaining axioms are non-standard.''} They are
registered, class-labeled, and deliberately visible; the program
nowhere claims they are proved. The reconciliation exists so that
no document can quietly re-count them.
\item \emph{``The 119{,}859-module batch tier poisons the
metrics.''} The caliber standard prohibits merging tiers: the
batch rate is reported separately (0\% honest rate by
construction), the one contaminated module is quarantined, and
the 30.2\% alarm turned out to be a caliber artifact --- which is
precisely why tiers and calibers are enforced in CI.
\item \emph{``Your own numbers disagreed; why believe the current
ones?''} The drift is the finding. Six documented conflicts
(C1--C6) were resolved through a named-caliber reconciliation
whose every step is reproducible; the errata are surface-visible,
never silent. A numbers system that can exhibit its own
corrections is the strongest thing an auditor can be offered.
\end{enumerate}

\subsection{Future work}
\label{sec:future}
The registered program, in order: (i) the metatheorem track ---
formalizing the conservativity of the three conversion families,
which would convert this paper's taxonomy from an audit practice
into a method with proved properties; (ii) evaluation of the
residual: the 239 primitive and 14 schema axioms need a grading
that distinguishes permanent physics postulates from redeemable
mathematical debt, and the Agda 131 require a machine that can
compile them; (iii) a non-delegation theorem target on the
Chern--Simons chain (the registered candidate: a weight-group
decomposition theorem generalizing T6's integer-sector
invariance), the first prospective THEOREM that would be neither
a wrapper nor a cast identity; (iv) containerized reproducibility
(Dockerfile, pinned toolchain, expected-output snapshots), closing
the P1-8 gap; (v) venue checkpoints per the journal-matching
matrix~\cite{sylvad2matrix2026}: ITP 2027 as the mid-way conference target, CPP 2028 as
the main line, JAR as the journal target this draft is shaped
for. The blind registry keeps its own schedule: $\alpha^{-1} =
137$ stays a frozen prediction until a first-principles
computation or its refutation arrives.

\subsection{Ethics and transparency}
\label{sec:ethics}
The library and this paper were produced with AI assistance under
human tasking and review. The repository's transparency conventions
(AI-assistance markers on every artifact, deleted-claims ledger,
no-silent-deletion rule, CI-enforced honesty markers) apply to this
manuscript: canonical numbers, bracketed baselines, draft status,
and pending work items are visible by construction, and the
supplementary material includes the build log of the exact
\LaTeX{} compile that produced the accompanying PDF. One decision
deserves explicit statement: the program does \emph{not} delete the
fabricated or vacuous material it finds --- it labels, quarantines,
or converts it, leaving the evidence in place. For an AI-assisted
corpus, an audit trail that survives is worth more than a clean
but unverifiable history.
